\documentclass[11pt]{article}
\usepackage[margin=1in]{geometry}
\usepackage{amsmath,amssymb,amsfonts,mathtools,bm}
\usepackage[T1]{fontenc}
\usepackage[utf8]{inputenc}
\title{Aligned Mathematical Structures}
\author{Source-backed grouped formula sample}
\date{}
\begin{document}
\maketitle
\section{Expressions}
\paragraph{Expression 1.} The following expression is taken from a source-backed formula corpus.

\begin{align*}F^{(L)}_{p;r}(q) =q^{-r(r-1)} \times\left\{ \begin{array}{ll}F^{(L)}_{p-2} \left[ \begin{array}{c}{\bf Q}_{r,1} \\ {\bf 0} \end{array} \right]({\bf e}_r | q) &\mbox {if $L \not\equiv r-1$ mod $2$}, \\F^{(L)}_{p-2} \left[ \begin{array}{c}{\bf Q}_{p-r,p} \\ {\bf 0} \end{array} \right]({\bf e}_{p-r} | q) &\text{if $L \equiv r-1$ mod $2$},\end{array}\right.\end{align*}

\paragraph{Expression 2.} The following expression is taken from a source-backed formula corpus.

\begin{align*}\begin{array}{c}u(\tau )=e^{\sqrt{\gamma ^2+m^2/2}\,\tau }\left[ u_0\cosh (\gamma \tau)-\left( u_0\sqrt{\gamma ^2+m^2/2}-\dot u_0\right) \frac 1\gamma \sinh(\gamma \tau )\right] \\v(\tau )=e^{-\sqrt{\gamma ^2+m^2/2}\,\tau }\left[ v_0\cosh (\gamma \tau)+\left( v_0\sqrt{\gamma ^2+m^2/2}+\dot v_0\right) \frac 1\gamma \sinh(\gamma \tau )\right] ,\end{array}\end{align*}

\paragraph{Expression 3.} The following expression is taken from a source-backed formula corpus.

\begin{align*}\begin{array}{l}\mathbf{E}(\mathbf{x},t) =\displaystyle\frac{1}{4\pi}\int\rho_e(\mathbf{x}',t-r)\frac{\mathbf{r}}{r^3}\,d^3x'\\~\\+\displaystyle\frac{1}{4\pi}\int\mathbf{J}_m(\mathbf{x}',t-r)\,{\bf\times}\,\frac{\mathbf{r}}{r^3}\,d^3x'-\displaystyle\frac{1}{4\pi}\int\frac{1}{r}\frac{\partial\mathbf{J}_m(\mathbf{x}',t-r)}{\partial t}\,d^3x' ,\end{array}\end{align*}

\paragraph{Expression 4.} The following expression is taken from a source-backed formula corpus.

\begin{align*}\Gamma ^{5}=\left( \begin{array}{cc}0 & 1_{2} \\ 1_{2} & 0\end{array}\right);\quad \Gamma ^{i}=\left( \begin{array}{cc}0 & \sigma ^{i} \\ -\sigma ^{i} & 0\end{array}\right);\quad \Gamma ^{0}=\left( \begin{array}{cc}1_{2} & 0 \\ 0 & -1_{2}\end{array}\right);\quad \Gamma ^{4}=-i\Gamma ^{5}=\left( \begin{array}{cc}0 & i_{2} \\ i_{2} & 0\end{array}\right).\end{align*}

\paragraph{Expression 5.} The following expression is taken from a source-backed formula corpus.

\begin{align*}(II) \hspace{3mm} \left(\begin{array}{c}x_1+{\rm i}y_1 \\x_2+{\rm i}y_2 \\x_3+{\rm i}y_3 \\x_4+{\rm i}y_4\end{array}\right) \hspace{3mm} \leftrightarrow \hspace{3mm}\frac{1}{\sqrt{2}}\left(\begin{array}{c}-(y_2+y_4) +{\rm j}(y_2-y_4) \\(y_1 + y_3) -{\rm j}(y_1-y_3) \\(x_2 + x_4) -{\rm j}(x_2 - x_4) \\-(x_1 + x_3) +{\rm j}(x_1 - x_3)\end{array}\right),\end{align*}
\end{document}
